\subsection{Unmixing Pairs Via Clustering}\label{sec:cluster}
We now show how to take flagged sets of columns from the tria runs and recover
pairs. Here, the flagged sets of columns mean the set $W$
to which sets of columns flagged as ``random" were added in
Algorithm~\ref{alg:main_analysis}
after each good trial.

The idea is simple once the flagged set of a single trial is looked
at the right way. In a trial shortening, a twin pair can be in only three
states. If both members are kept, then both are dirty columns, so
both are flagged with very very high probability.
If exactly one member is kept, and then the survivor
is derandomized to a clean column per the analysis of the Couvrier et al
paper~\cite{CLT18},
so that column is not flagged, and the
other member is absent from the trial altogether. The final case is where both
members
are hit by the shortening, and neither appears. In all three cases
it is with very very high probability the case that either both are flagged or neither
are flageed.

Finally, a column that belongs to no twin pair, (a pure GRS column), is never
flagged at all.

As a result, if both members of a pair are correctly flagged with high enough
probability,
then the list of trials in
which column A was flagged  is the same as the list of columns for which its
twin column B was flagged.
Meanwhile, it is an empty list for a pure column, and for other twin pairs, the
list is non-empty but with
high probability over the random choices of shortenings distinct. Pair
identification thus becomes grouping
columns by their \emph{pattern}.

\subsubsection{The algorithm}

For each original public column $i$, record its pattern
$b_i \in \{0, 1\}^T$, where $b_i(t) = 1$ if $i$ was flagged in
trial $t$ and $0$ otherwise. Then proceed as follows.

\begin{algorithm}[H]
    \caption{Flag-pattern clustering}
    \label{alg:cluster}
    \KwIn{Flagged sets $V_{1}, \dots, V_{T}$ of $T$ trials,
        mapped back
        to original public-column indices}
    \KwOut{Sets of twin pairs and pure columns; anomalies, if any}
    For each column $i$, form its pattern $b_i$, with $b_i(t) = 1$
    iff $i \in V_t$\;
    Group columns by \emph{exact equality} of patterns\;
    A group of size $2$ with a nonzero pattern is a twin pair\;
    The group with the all-zero pattern is the pure columns\;
    A group of size $1$ with a nonzero pattern is an orphan\;
    A group of size $3$ or more is oversized, and is reported, not
    paired\;
    Compare the size of the all-zero group with the expected
    number $n - w$ of pure columns, and report any excess\;
\end{algorithm}



\subsubsection{Trials for Success}\label{sec:cluster-planning}

We now compute how many trials (assuming that pairs are with overwhelming probability correctly flagged together) 
are necessary to identify and unmix pairs with high probability. Fix a pair. 
The probability both columns of the pair appear in the shortening is \[
    p \;=\; \frac{(n')(n'-1)}{(n+w)(n+w-1)},
\]
the probability that a given pair survives a random shortening to
$n'$ out of $N$ columns. Two distinct pairs separate in a recorded trial
except with probability at
most $p^2 + (1 - p)^2$, when either both pairs appear in the shortening, or neither pair appears in the shortening.


Note that we heuristically treat the pairs as independent here. A union bound over the $\binom{w}{2}$ pairs of
pairs then says that $T$ trials separate every pair from every
other pair with failure probability at most $\epsilon$ as soon as
\[
    \binom{w}{2} \bigl(p^2 + (1 - p)^2\bigr)^T \;\leq\; \epsilon,
\]

and we take $T_{\mathrm{sep}}$ to be the smallest such $T$. A pair
is invisible, and so hides in the pure set, if it is never unhit
in $T$ trials, which happens with probability $(1 - p)^T$ per pair;
a union bound over the $w$ pairs gives the visibility requirement
\[
    w (1 - p)^T \;\leq\; 0.01,
\]
with $T_{\mathrm{vis}}$ the smallest such $T$. Clustering runs $
    T \;=\; \max(T_{\mathrm{sep}}, T_{\mathrm{vis}})
$.

\subsubsection{Empirical Validation}\label{sec:cluster-validation}

Clustering was validated at several points, where the truth
was known from the key and every trial could be checked against
it. 

These parameters include the actual ID-0 parameters. 

\begin{table}[htbp]
    \centering
    \footnotesize
    \begin{tabular}{lcccc}
        \toprule
        Scheme Parameters ($q, n, k, w$) & Attack Parameters ($n',(s,d)$)  & Trials & Flagging exact & Pairs recovered \\
        \midrule
        $(256, 198, 165, 33)$, $(90, 2, 3)$           & $23/23$        & $33/33$    \\
        $(256, 119, 79, 31)$, $(101, 2, 3)$  & $25/25$        & $31/31$   \\
        $(1024, 630, 470, 160)$, $(366, 2, 3)$ & $41/41$ & $160/160$ \\ 
        \bottomrule
    \end{tabular}
    \caption{Cluster validation at reduced parameters. ``Flagging
        exact'' means the flagged set equalled the unhit pairs' columns
        in that many trials.}
    \label{tab:cluster-validation}
\end{table}


\paragraph{ID-0, end to end.} The shipped ID-0 set has since been
run end to end with cluster Phase~1, on a fresh key generated with
the NIST DRBG key generation of the implementation, at the adopted
point $(n', k', (s, d)) = (366, 46, (2, 3))$, with $41$ trials on
$12$ threads. The run succeeded completely: cluster aggregation
over the $41$ trials returned $160$ pairs, $470$ pure columns,
$0$ orphans, and $0$ oversized groups; Phase~2 unmixed $160/160$
pairs, Sidelnikov--Shestakov recovered the correct GRS code, and
verification passed with $160/160$ non-degenerate pairs correct
and $0$ false positives. The run took a little over $3$ hours of
wall time across those $12$ threads to break ID-0. That is the
Phase~1 trial time dominating the run: $41$ trials in batches of
$12$, each trial a single reuse elimination, with aggregation,
Phase~2 (SS bootstrap $0.25$\,s), and verification taking minutes
by comparison. ID-1 and ID-2 have not been run end to end with
cluster Phase~1; their figures below remain model estimates.
